\documentclass[12pt]{article}
\usepackage{amsmath,amssymb,amsfonts,amsthm}
\usepackage{braket}
\usepackage{hyperref}
\usepackage{geometry}
\usepackage{authblk}
\geometry{margin=1in}

\newtheorem{theorem}{Theorem}[section]
\newtheorem{lemma}[theorem]{Lemma}
\newtheorem{proposition}[theorem]{Proposition}
\newtheorem{corollary}[theorem]{Corollary}
\theoremstyle{definition}
\newtheorem{definition}[theorem]{Definition}
\newtheorem{assumption}[theorem]{Assumption}
\theoremstyle{remark}
\newtheorem{remark}[theorem]{Remark}

\newcommand{\Hthree}{\mathbb{H}^3}
\newcommand{\Sphere}{\mathbb{S}^2}
\newcommand{\GammaG}{\Gamma}
\newcommand{\U}{\mathrm{U}}
\newcommand{\calH}{\mathcal{H}}
\newcommand{\bdry}{\partial\mathbb{H}^3}
\newcommand{\calHS}{\mathcal{H}_S}
\newcommand{\calHT}{\mathcal{H}_T}

\title{Discrete Mixing Operators from Boundary Sector Geometry}
\date{}
\author[1]{Marvin L.\ Gentry}
\affil[1]{Independent Researcher\\
  \texttt{drmlgentry@protonmail.com}\\
  ORCID: 0009-0006-4550-2663}

\begin{document}
\maketitle

\begin{abstract}
We develop a kinematic framework in which unitary mixing matrices arise
as overlap operators between sector subspaces generated by boundary-localized
states in a Hilbert space associated with a compact hyperbolic $3$-manifold.
Assuming a discrete assignment of sectors to boundary directions,
we prove that the resulting mixing matrices are uniquely determined,
up to sector-local rephasings, by the underlying boundary configuration.
Our main result (Theorem~\ref{thm:discrete-classes}) establishes that the
set of admissible mixing matrices is discrete in $\U(n)$ modulo
$\U(1)^{2n}$, and that no continuous deformation of parameters can
produce continuous drift in any rephasing-invariant mixing observable.
Rephasing-invariant quantities depend only on the discrete boundary data
and are stable under bounded deformations preserving the assignment rule.
The construction is purely geometric and kinematic, requiring no dynamical
input; it provides a model-independent structural constraint on theories
in which fermions are associated with boundary directions in a compactified
hyperbolic geometry.
\end{abstract}

%--------------------------------------------------
\section{Introduction}
%--------------------------------------------------

Mixing matrices, such as the CKM matrix in the quark sector and the PMNS
matrix in the lepton sector, encode the mismatch between flavor and mass
eigenstates in the Standard Model~\cite{PDGreview}.  Despite their
fundamental importance, their origin remains unexplained by the gauge
structure alone.  Many attempts to derive these matrices from first
principles invoke discrete flavor symmetries, extra dimensions, or
string-inspired constructions~\cite{AltarelliFeruglio,KingLuhn}.  In all
such approaches, the mixing matrix ultimately measures the overlap between
states associated with different sectors of the theory.

In this paper we take a kinematic, model-independent perspective.
We show that if fermions are assigned to distinct boundary directions of a
compact hyperbolic $3$-manifold, the mixing matrices that arise from the
overlap of the corresponding boundary-localized states are determined solely
by the discrete boundary data.  Moreover, these matrices form a discrete
set modulo sector-local rephasings---a strong structural constraint that
any realistic model must satisfy.  A corollary (Corollary~\ref{cor:nodrift})
is that no continuous deformation of physical parameters can produce
continuous drift in any rephasing-invariant mixing observable.

The construction is deliberately minimal: a Hilbert space of boundary data,
a family of boundary-localized states, and a discrete assignment of fermions
to boundary directions.  No dynamics, no Lagrangian, and no ultraviolet
completion are assumed.  All results follow from elementary Hilbert space
geometry and the discreteness of the boundary labels.

The paper is organized as follows.  Section~\ref{sec:geom} establishes
the geometric setup and the discreteness assumption.
Section~\ref{sec:states} introduces the boundary Hilbert space and
localized states.  Sector subspaces and preferred bases are defined in
Section~\ref{sec:sectors}.  Mixing matrices appear as overlap operators in
Section~\ref{sec:mixing}.  Rephasing invariance and its consequences are
discussed in Section~\ref{sec:rephasing}.  Section~\ref{sec:stability}
establishes stability under bounded deformations and the main discreteness
theorem.  We conclude in Section~\ref{sec:conclusion}.

The geometric construction developed here provides the theoretical
foundation for the companion results in~\cite{Gentry:CKM,Gentry:PMNS},
where the CKM and PMNS mixing matrices are derived explicitly from the
$K$- and $N$-factors of the Iwasawa decomposition of compact hyperbolic
3-manifold holonomy~\cite{Gentry:CKM,Gentry:PMNS}.

%--------------------------------------------------
\section{Geometric Setup and the Discreteness Assumption}
\label{sec:geom}
%--------------------------------------------------

Let $\GammaG$ be a torsion-free discrete group acting properly discontinuously
and cocompactly on $\Hthree$ by orientation-preserving isometries.
The quotient
\[
  M = \Hthree/\GammaG
\]
is a compact oriented hyperbolic $3$-manifold.
Its boundary at infinity is the sphere $\Sphere = \bdry$.

\begin{definition}[Boundary direction]\label{def:bdrydir}
A \emph{boundary direction} is a point $\xi\in\Sphere$.
A group element $\gamma\in\GammaG$ is \emph{associated with} $\xi$ when
the geodesic ray from a fixed basepoint $o\in\Hthree$ through $\gamma\cdot o$
converges to~$\xi$.
\end{definition}

\begin{assumption}[Discrete sector assignment]\label{ass:discrete}
Each fermion flavor $f\in F$ is assigned a unique boundary direction
$\xi(f)\in\Sphere$ by a deterministic, discrete selection rule:
the image $\Xi=\{\xi(f):f\in F\}$ is finite, and the assignment is
locally rigid.
A concrete realization uses the ideal endpoint of the geodesic ray to
the minimal representative $\gamma_f\cdot o$; uniqueness and rigidity
follow from~\cite{bridsonhaefliger}.
\end{assumption}

%--------------------------------------------------
\section{Boundary Hilbert Space and Localized States}
\label{sec:states}
%--------------------------------------------------

Let $\calH$ be a separable complex Hilbert space carrying boundary data.
A concrete realization is $L^2(\Sphere)$, but any separable Hilbert space
admitting a continuous family of normalized vectors $\ket{\xi}$ suffices.

\begin{definition}[Boundary-localized state]\label{def:localized}
For each $\xi\in\Sphere$, fix a normalized vector $\ket{\xi}\in\calH$
such that the map $\xi\mapsto\ket{\xi}$ is continuous and
$\braket{\xi|\zeta}=1$ if and only if $\xi=\zeta$.
\end{definition}

The continuity condition ensures that nearby boundary directions give
nearly parallel states; the normalization condition implies that distinct
directions produce linearly independent vectors.

%--------------------------------------------------
\section{Sector Subspaces and Preferred Bases}
\label{sec:sectors}
%--------------------------------------------------

Let $F$ be a finite set of fermion labels.

\begin{definition}[Sector]\label{def:sector}
A \emph{sector} $S$ is a finite ordered subset of $F$.
The ordering is fixed once and for all (e.g., by increasing mass).
\end{definition}

\begin{definition}[Sector subspace]\label{def:sectorspan}
The \emph{sector subspace} of $S$ is
\[
  \calHS = \operatorname{span}\{\ket{\xi(f)}:f\in S\}\subset\calH.
\]
Linear independence of the generating set implies $\dim\calHS=|S|$.
\end{definition}

\begin{definition}[Preferred orthonormal basis]\label{def:preferred}
The \emph{preferred ONB} of $\calHS$ is the Gram--Schmidt
orthonormalization of $(\ket{\xi(f_1)},\ldots,\ket{\xi(f_{n_S})})$
in the fixed sector ordering, with all pivot coefficients real and positive.
The resulting basis is denoted $\{\ket{S,1},\ldots,\ket{S,n}\}$.
\end{definition}

The preferred basis depends on the sector ordering, but different orderings
yield unitarily equivalent mixing matrices up to permutations, which can
be absorbed into the rephasing freedom.

%--------------------------------------------------
\section{Mixing Matrices as Overlap Operators}
\label{sec:mixing}
%--------------------------------------------------

Let $S$ and $T$ be two sectors of equal cardinality $n$.

\begin{definition}[Mixing matrix]\label{def:mixing}
The \emph{mixing matrix} $U_{ST}$ is the $n\times n$ matrix with entries
\[
  (U_{ST})_{ji} = \braket{T,j|S,i}.
\]
\end{definition}

\begin{proposition}[Unitarity]\label{prop:unitary}
If $\calHS=\calHT$, then $U_{ST}\in\U(n)$.
In general, $U_{ST}$ is the matrix of the orthogonal projection
$P_T|_{\calHS}$ in the preferred bases and is a partial isometry.
\end{proposition}

\begin{proof}
When the subspaces coincide, the two preferred bases are orthonormal bases
of the same space, so the change-of-basis matrix is unitary.  In the
general case, the operator $P_T|_{\calHS}$ is a partial isometry, and its
matrix in the preferred bases is exactly $U_{ST}$.
\end{proof}

\begin{theorem}[Geometric dependence]\label{thm:geom}
For fixed ordering conventions, $U_{ST}$ is determined entirely by the
finite boundary configuration
\[
  \Xi_{ST} = \{\xi(f):f\in S\cup T\}
\]
and the inner product on $\calH$.
\end{theorem}

\begin{proof}
All Gram--Schmidt coefficients are algebraic expressions in inner products
$\braket{\xi(f)|\xi(g)}$, which depend only on the points in $\Xi_{ST}$
and the fixed inner product.
\end{proof}

%--------------------------------------------------
\section{Rephasing Invariance}
\label{sec:rephasing}
%--------------------------------------------------

In physical applications, unphysical phases can be removed by redefining
the phases of fermion fields.  Within a given sector, such redefinitions
correspond to multiplying each basis vector $\ket{S,i}$ by an independent
phase $e^{i\alpha_i}$, inducing
\[
  U_{ST} \mapsto D_T\,U_{ST}\,D_S^\dagger,
\]
where $D_S,D_T\in\U(1)^n$ are diagonal unitary matrices.
Physical information is therefore contained in the equivalence class of
$U_{ST}$ in $\U(n)/(\U(1)^n\times\U(1)^n)$.

\begin{proposition}[Rephasing-invariant quantities]\label{prop:rephinv}
The following are invariant under sector-local rephasings:
\begin{enumerate}
\item[\emph{(i)}] the moduli $|(U_{ST})_{ji}|^2$;
\item[\emph{(ii)}] the quartic cyclic products
  $J_{ji,j'i'}=(U_{ST})_{ji}(U_{ST})_{j'i'}
   \overline{(U_{ST})_{ji'}}\,\overline{(U_{ST})_{j'i}}$;
\item[\emph{(iii)}] $\det(U_{ST})$ up to an overall phase.
\end{enumerate}
\end{proposition}

The quantities in (ii) generalize the Jarlskog invariant~\cite{Jarlskog}
to the present geometric setting.

%--------------------------------------------------
\section{Stability and Discreteness of Mixing Classes}
\label{sec:stability}
%--------------------------------------------------

\begin{definition}[Bounded deformation]
Let $d_{\Sphere}$ be a metric on $\Sphere$ compatible with its topology.
A \emph{deformation of magnitude $\varepsilon$} replaces each $\xi(f)$
by a point $\xi'(f)$ with $d_{\Sphere}(\xi(f),\xi'(f))\leq\varepsilon$.
\end{definition}

\begin{theorem}[Stability]\label{thm:stability}
Let $\delta_{\min}=\min_{f\neq f'}d_{\Sphere}(\xi(f),\xi(f'))$ be the
minimal pairwise separation of the original boundary directions, and set
$\varepsilon_0=\delta_{\min}/2$.
For any deformation of magnitude $\varepsilon\leq\varepsilon_0$, the
perturbed boundary configuration $\Xi'$ determines the same preferred
bases and the same mixing matrix $U_{ST}$.
\end{theorem}

\begin{proof}
For $\varepsilon\leq\varepsilon_0$, each perturbed point $\xi'(f)$ remains
closer to $\xi(f)$ than to any other $\xi(g)$ with $g\neq f$.
In particular, the $\xi'(f)$ remain distinct, so linear independence is
preserved.  Since Gram--Schmidt coefficients depend continuously on inner
products, and since $\braket{\xi'(f)|\xi'(g)}$ varies continuously with
$\varepsilon$, the preferred ONBs vary continuously.
However, the discrete assignment rule (Assumption~\ref{ass:discrete}) is
unchanged: each $\xi'(f)$ belongs to the same assignment class as $\xi(f)$.
By Theorem~\ref{thm:geom}, the mixing matrix is a function of the
configuration alone, and since the configuration class is unchanged,
$U_{ST}$ is unchanged.
\end{proof}

\begin{theorem}[Discreteness of mixing classes]\label{thm:discrete-classes}
Under Assumption~\ref{ass:discrete}, the set of accessible mixing matrices
\[
  \mathcal{M}_{ST} = \{U_{ST} : \Xi_{ST}\ \text{admissible}\}
\]
is discrete in $\U(n)/(\U(1)^n\times\U(1)^n)$.
\end{theorem}

\begin{proof}
The set of admissible configurations $\Xi_{ST}$ is finite.
By Theorem~\ref{thm:geom}, the map $\Xi_{ST}\mapsto U_{ST}$ is continuous,
hence maps a finite set to a finite set.
Quotienting by the rephasing action preserves finiteness, so
$\mathcal{M}_{ST}$ is a finite---hence discrete---subset of the quotient.
\end{proof}

\begin{corollary}[No continuous drift]\label{cor:nodrift}
No continuous deformation of physical parameters can produce continuous
drift in any rephasing-invariant mixing observable.
\end{corollary}

\begin{proof}
Rephasing-invariant observables are continuous functions on
$\U(n)/(\U(1)^n\times\U(1)^n)$.  Their restriction to the discrete set
$\mathcal{M}_{ST}$ takes only finitely many values; continuous drift is
therefore impossible.
\end{proof}

%--------------------------------------------------
\section{Conclusion}
\label{sec:conclusion}
%--------------------------------------------------

We have presented a purely kinematic construction of mixing matrices from
boundary-localized states in a compact hyperbolic $3$-manifold.
The only inputs are a Hilbert space, a continuous family of normalized
states $\ket{\xi}$, and a discrete assignment of fermions to boundary
directions $\xi(f)$.
From these, each pair of sectors gives rise to a unitary (or partially
isometric) mixing matrix, uniquely determined up to sector-local rephasings.
Our main theorem establishes that under the discreteness assumption the set
of admissible mixing matrices is finite modulo rephasings---a strong
structural constraint on any geometric theory of flavor mixing.
A direct corollary is the absence of continuous drift in all
rephasing-invariant observables.

This result is model-independent: it holds regardless of the dynamics that
might generate the boundary assignment.  In a companion paper we show how
boundary directions can arise from unique minimal representatives in a
cocompact group action~\cite{bridsonhaefliger,ratcliffe},
linking the present kinematic framework to concrete arithmetic constructions.
Future work could investigate whether the observed CKM and PMNS matrices
belong to the discrete set for natural choices of boundary directions, or
whether additional geometric structure is required~\cite{PDGreview}.

\bibliographystyle{plain}
\bibliography{gentry-flavor-mixing}

\end{document}